\section{Evaluación y modos de falla}

\subsection{Jerarquía de métricas y gobernanza}
\begin{table}[H]
\centering
\begin{tabular}{p{3.5cm}p{4cm}p{3.5cm}p{3cm}}
\toprule
Métrica & Descripción & Frecuencia de recolección & Condición de activación \\
\midrule
Proof success rate & Si Lean demuestra con éxito la proposición & real-time per episode & success < 2\% → ampliar search \\
Entropy drift & Policy entropy drop & epoch & drop > 30\% → drift gate log \\
Chunk hit rate & Tasa de aciertos de la cache del search track & per batch & < 80\% → adjust overlap \\
Human edit count & Número de traces que requieren intervención humana & weekly & > 5/1000 → governance review \\
\bottomrule
\end{tabular}
\caption{Métricas de evaluación multinivel y estrategias de respuesta de AlphaProof}
\end{table}

Hubert mencionó que las diapositivas 17-19 muestran el dashboard completo de métricas: la parte superior corresponde a proof success/distribution, el lado izquierdo al entropy drift heatmap y el lado derecho al user feedback count. El dashboard está conectado directamente con el alert pipeline del governance board, y en cuanto una track cruza el threshold, automáticamente envía el replay log al reviewer.

\begin{importantbox}{El ciclo cerrado de evaluación y gobernanza}
1) Recolección paralela de múltiples métricas; 2) cuando el proof success rate cae por debajo del baseline, se aumenta el search depth y se registra el replay; 3) el entropy drift y los manual edits determinan en conjunto si se necesita una governance review.
\end{importantbox}

\subsection{Análisis de los modos de falla}
Las fallas registradas por AlphaProof se concentran sobre todo en los siguientes patterns: A) semantic gap: error de formalizer translation; B) search drift: la policy entra en una low-entropy zone; C) Lean constraints: un type mismatch provoca el colapso del goal tree; D) compute budget exceed. Hubert lo representa de manera gráfica con el «failure funnel»: la mayoría de las failure se filtran en la etapa de initial seeding, y solo una minoría hace drift via chunk gating hasta llegar al governance board.

\begin{warningbox}{Advertencia sobre fallas comunes}
El semantic gap requiere volver al formalizer; una vez que se activa el drift gate, primero se analiza el entropy heatmap; los errores de tipo requieren revisar el Lean goal stack; cuando ocurre un budget exceed, se baja directamente al low-latency track.
\end{warningbox}

\subsection{Reproducción y robustez}
Cada failure trace se escribe en `failure-repro.json`, e incluye time stamp, search track, entropy y human comment. El equipo sube estos trace al experiment repo para generar el `failure curriculum`, que permite que el RL agent haga un nuevo warm start. Este proceso está supervisado por el governance board, que garantiza que cada failure log pueda convertirse en el next training seed.

\begin{knowledgebox}{Elementos del failure curriculum}
1) Trace + entropy + comment conforman una entry; 2) la entry entra en el Lemma buffer; 3) el training schedule la usa como positive/negative sample; 4) Result logged back to dashboard.
\end{knowledgebox}

\subsection{Resumen de la sección}
La capa de evaluación y modos de falla conecta la percepción (métricas), el diagnóstico (failure taxonomy) y la reproducción (failure curriculum), lo que permite que AlphaProof, a lo largo de sus iteraciones, convierta las fallas en datos de entrenamiento controlables.
\newpage

